problem:  
Let \((M^m,g_M)\) be a complete noncompact Riemannian manifold with \(\mathrm{Ric}_M\ge0\) and supporting the sharp Euclidean–type Sobolev inequality  
\[
\|v\|_{L^{2m/(m-2)}}^2 \le S_M\,\|\nabla v\|_{L^2}^2, \quad v\in C_c^\infty(M).
\]

Let \((N^n,g_N)\) be a simply connected Cartan–Hadamard manifold with **anisotropic curvature decay** in polar coordinates about some reference point \(p\in N\): there exist constants \(C_1,C_2>0\) such that for all sufficiently large \(r=d_N(p,x)\),
\[
\begin{aligned}
-\,C_2(1+r)^{-\alpha_t} &\le K_N(e_t,e_s) \le -\,C_1(1+r)^{-\alpha_t},\\
-\,C_2(1+r)^{-\alpha_r} &\le K_N(e_r,e_t) \le -\,C_1(1+r)^{-\alpha_r},
\end{aligned}
\]
where \(e_r\) denotes the radial unit vector and \(e_t,e_s\) are orthonormal tangential directions.

Suppose \(u:M\to N\) is a harmonic map of finite energy
\[
E(u) = \frac{1}{2}\int_M |\nabla u|^2 < \infty.
\]

---

**Conjectural Problem:**  
There exists a critical *anisotropic rigidity threshold* \((\alpha_r^*,\alpha_t^*)\) depending only on \(m\) such that:

1. If \(\min\{\alpha_r,\alpha_t\} < \alpha_r^*=\alpha_t^* = 2,\) then every finite–energy harmonic map \(u:M\to N\) is constant.  
2. If \(\min\{\alpha_r,\alpha_t\} > 2\), there exist examples of nonconstant finite–energy harmonic maps for some \((M,N)\) satisfying the above decay laws.  
3. (Anisotropic question) If \(\alpha_r < 2 < \alpha_t\) (or conversely), does rigidity persist or fail?  Equivalently, does *faster tangential flattening* compensate for *slower radial decay*, or is the isotropic quadratic rate universally critical?

---

Why is it a "good" problem:  
This refined formulation is mathematically coherent and targets a subtle, concrete conjecture: whether **directionally dependent curvature decay** affects the sharp quadratic rigidity threshold for harmonic maps. It isolates a clear analytic‑geometric quantity (the pair \((\alpha_r,\alpha_t)\)) and asks for the precise balance between curvature components that maintains or breaks vanishing theorems.  
Resolving it would deepen understanding of how anisotropy in a manifold’s asymptotic geometry influences Bochner‑type energy estimates, bridging sectional curvature geometry with PDE anisotropy. Its statement is simple—two decay exponents—but it probes a frontier mechanism linking curvature tensors and harmonic‑map rigidity, embodying the “narrow entrance, profound depth’’ ideal of a good research problem.